\scanpage{30}\mspage{29}
\uline{Cont 4} If $F,G,$ \deleted{$H$ are figures}\inserted{are two} figures, such that every multistratum of $F$ is compatible with every multistratum of $G^{*}$, then $F$ and $G$ are compatible (the converse being obvious)

% slanted
\margin{\uline{NB} Then $F$ and $G$ are set-theoretically compatible.}

\deleted{\uline{Corollary} If $F,G,H$ are \unsure{. all|. three} $(\unsure{F_i|F_j}]$}

\deleted{Équi} This implies

\uline{Corollary} Let $(F_i)_{i\in I}$ be a family \deleted{\unsure{finite}} of figures. For \deleted{the $F_i$ to be} this family to be a configuration, it is necessary and sufficient that for every $i\ne j$ and every multistratum $\widehat{X}\deleted{\in}$ of $F_i$, $\widehat{Y}$ of $F_j$, $\widehat{X}$ and $\widehat{Y}$ be compatible. If $I$ is finite, this implies that $F=\bigcup F_i$ is a figure, and \inserted{the} $F_i$ \deleted{is} are subfigures.

\deleted{\uline{NB} Conversely, the validity of this statement for \deleted{some} every pair $\{F,G\}$ of figures (namely, \deleted{if the} the criterion of}

\uline{Countries \deleted{set-theoretic}\inserted{of points}} The country is called ``set-theoretic'', or ``\uline{country of points}'' (\unsure{that is|here|there|i.e.} its loci take the name of \uline{points}) if two set-theoretically compatible figures are compatible (which thus makes Cont 4 tautological)
